\chapter{Rendering}

\section{What a Renderer Actually Does}

Chapter 25 established that a renderer is a client of the world database, not a stage of the core write loop, and Chapter 9 established that the observation operator \(O_\theta\) actually has two directions: forward, reading committed appearance into a perceivable output, and inverse, proposing a candidate appearance for whatever a query touches that has not yet been resolved. Both chapters deferred the concrete implementation of either direction. This chapter supplies it, and the central complication this chapter has to resolve is that the two directions, despite both being invoked by what looks from the outside like a single act of rendering, have entirely different obligations: one is a read, and the other, however it may appear to a user simply looking at a scene, is a write.

\section{The Camera as a Query}

A camera, or more generally any observer's request for a view of the world, is not a passive framebuffer waiting to be filled. It is a query, issued against the world database's interface from Chapter 26: a region of the domain, a requested resolution, and a specification of which of the five components the request actually needs. This last part matters more than it might first appear. Chapter 7 already established that different observation operators recover different partial slices of \(M_t\), a camera reading mostly z, a diagnostic tool reading Φ or R instead, and this chapter's architecture takes that claim literally: a visual renderer's query asks the world database for appearance at fine spatial resolution and coherence or entropy only coarsely, if at all, while a structural inspection tool's query might ask for Φ and R at fine resolution and skip appearance almost entirely. There is no single, universal rendering query. There is a family of queries, each shaped by what the requesting observer actually needs to see.

\section{Decoding a Resolved Region}

For a region the query touches that is already resolved, low entropy, committed appearance, decoding is the straightforward half of this chapter's work, and it is worth being honest that this book has little to add to it beyond situating it correctly. A decoder, whatever neural or classical machinery implements it, takes the stored representation of z at the queried region and transforms it into whatever output modality the observer requires, pixels for a visual renderer, waveform samples for an audio one, natural-language description for a text-based observer. This is well-trodden territory, and this book's contribution is not a new decoding technique but the guarantee everything upstream of the decoder now provides: the z being decoded is not freshly sampled for this particular query, the way a generative model with no persistent state would produce it, but retrieved, unchanged since it was last committed, from the world database itself. The decoder's job is translation, not invention, whenever the region it is asked to decode has already been resolved.

\section{Querying an Unresolved Region}

The harder and more consequential case is a query that touches a region still at high entropy, still represented only by \(L_t(x)\), Chapter 8's live possibility set, with no committed z to decode. A decoder cannot simply decode nothing; some output has to be returned to satisfy the query. It is worth being precise about what happens here, because getting this wrong reintroduces exactly the failure this entire book began by diagnosing.

The correct behavior is not for the renderer to generate a plausible appearance and hand it back directly, bypassing everything Chapters 12 through 15 established about how the field is permitted to change. Instead, the query triggers the inverse direction of \(O_\theta\): a candidate ẑ is generated, exactly as Chapter 9 described, and this candidate is submitted as a genuine proposal, entering the same proposal-projection-commit pipeline any other \(\Delta\psi\) would enter, per Chapter 25's architecture. The rendering request does not complete until this proposal has been projected against C and, assuming it survives, committed to the world database. Only then is the newly committed z decoded and returned to the observer. This is a slower path than decoding an already-resolved region, and it is slower for a principled reason: rendering an unresolved region is not a read at all, despite appearances. It is a write, disciplined by the same admissibility checking every other write in this book's architecture is subject to, and any implementation that shortcuts this path for the sake of rendering performance has, whether or not it recognizes the fact, reintroduced the wall problem this book opened with.

\section{Level of Detail as Honest Uncertainty}

Conventional rendering systems vary detail with distance from the observer, a familiar and well-understood optimization. This architecture has a second, independent axis to vary detail along, and it is worth naming because it is genuinely different in kind from the first. A region can be rendered at reduced fidelity not because it is far from the camera, but because it remains genuinely high in entropy, only partially resolved, and asking the decoder to produce more detail than the underlying \(L_t(x)\) actually supports would mean fabricating specificity the field itself does not yet have. This is not a performance shortcut. It is the rendering system being honest about what it knows, the visual equivalent of Chapter 8's insistence that entropy be allowed to remain wide wherever it has not been earned down. A distant, fully resolved wall can be rendered in coarse detail purely for performance reasons and still be, if the camera moved closer, revealed to hold a specific, already-committed crack. A nearby but genuinely unobserved region rendered in coarse detail is coarse for a different reason entirely: because there is, honestly, not yet more to show.

\section{Multiple Observers, One World}

Because most rendering queries against already-resolved regions are pure reads, many observers, human players, other agents, diagnostic tools, can query the same world database concurrently without any coordination between them, exactly as Chapter 25's client model intended. This changes the moment a query triggers section 27.4's invention path. An invention-triggering query is a proposal, and proposals, per Chapter 15, are subject to the same commutative and non-commutative arbitration any other simultaneous write is subject to. Two observers querying the same unresolved region at nearly the same moment are, functionally, two competing proposals for what that region should become, and the world database resolves this the same way it resolves any other conflict Chapter 15 described, not through any special-cased rendering logic. This book will not develop the full multi-observer treatment here; Chapter 29 takes it up directly, extending exactly this picture to the case where many observers, and many proposal-generating agents, are active in the same world at once.

\section{Looking Ahead}

This chapter has given the read side of this book's architecture its full engineering treatment: a camera as a query, a decoder for resolved content, and an inversion into the proposal pipeline for unresolved content, with detail honestly scaled to genuine uncertainty rather than merely to distance. The write side has, so far, only been engineered from the perspective of a query accidentally triggering it. Chapter 28 addresses writing on its own terms, as something an agent deliberately does rather than something a renderer happens to cause, building the interaction layer that turns an agent's intent into the \(\Delta\psi\) Chapter 13 already knew how to receive.
